\documentclass[11pt]{article}
% Shared preamble for the printable mini lessons (1.4 onward).
%   \documentclass[11pt]{article}
%   % Shared preamble for the printable mini lessons (1.4 onward).
%   \documentclass[11pt]{article}
%   % Shared preamble for the printable mini lessons (1.4 onward).
%   \documentclass[11pt]{article}
%   \input{../preamble}
% Build the key with:  pdflatex -jobname=<name>_KEY "\def\iskey{}\input{<name>}"
% Figures shared with the web page live in figures/Unit_X/<name>.tex and are
% pulled in with \fig{Unit_X/<name>}; the standalone package drops their own
% preamble, so their macros come from figures/preamble.tex, loaded here.
\input{../../figures/preamble}
\usepackage[margin=0.75in]{geometry}
\usepackage{array}
\usepackage{enumitem}
\usepackage{amsmath}
\usepackage{cancel}
\usepackage[helvet]{sfmath}  % math in the same sans font as the text
\usepackage{standalone}
\renewcommand{\ttdefault}{pcr}  % Courier: a real Type 1 typewriter font
\pagestyle{empty}
\setlength{\parindent}{0pt}
\setlist[enumerate,1]{leftmargin=1.5em,itemsep=1em}
\setlist[enumerate,2]{label=(\alph*),leftmargin=2em,itemsep=2.2em}

\newif\ifkey
\ifdefined\iskey \keytrue \fi
\colorlet{keyred}{red!75!black}
\newcommand{\ans}[1]{\ifkey\leavevmode{\color{keyred}\let\frac\dfrac #1}\fi}  % full-size fractions in answers
% vertical work space on the student copy, less on the key
\newcommand{\work}[1]{\ifkey\vspace{0.6em}\else\vspace{#1}\fi}
\newcommand{\fig}[1]{\input{../../figures/#1}}
% Title line: \lessontitle{Quartiles, IQR, and Box Plots}
% A long title gets the Name line on its own line above it.
\newlength{\titlewidth}
\newcommand{\lessontitle}[1]{%
  \settowidth{\titlewidth}{\Large\bfseries #1\ifkey\ KEY\fi}%
  \ifdim\titlewidth>4.3in
    \hfill Name \underline{\hspace{2in}}\par\vspace{0.5em}
    {\Large\bfseries #1\ifkey\ \textcolor{keyred}{KEY}\fi}\par\vspace{0.6em}
  \else
    {\Large\bfseries #1\ifkey\ \textcolor{keyred}{KEY}\fi}\hfill Name \underline{\hspace{2in}}\par\vspace{0.6em}
  \fi}
\newcommand{\partner}{\par\vspace{0.4em}Work with a partner. Compare answers when you finish each problem.\par}

% ---------- shared pieces added for AI and You ----------
% no hyphenated words anywhere: students read in a second language
\hyphenpenalty=10000 \exhyphenpenalty=10000
\usepackage{pifont}   % \ding{34} scissors
\newcommand{\blank}[1]{\underline{\hspace{#1}}}
% ragged table columns: no hyphenated words for students reading in a second language
\newcolumntype{R}[1]{>{\raggedright\arraybackslash}p{#1}}
% An answer written on a blank line: \ansline{width}{answer}
\newcommand{\ansline}[2]{\ifkey\underline{\makebox[#1][l]{\,\textcolor{keyred}{#2}}}\else\blank{#1}\fi}
% Unit 1 class timeline: a card to cut out and post on the wall.
%   \timelinecard{1950}{Alan Turing asks ...}
\newcommand{\timelinecard}[3][8.3]{%
  \begin{tikzpicture}
    \pgfmathsetmacro{\tw}{#1-0.8}
    \draw[dashed, line width=0.8pt, rounded corners=3pt] (0,0) rectangle (#1,4.7);
    \node[anchor=north west, font=\huge\bfseries] at (0.3,4.45) {#2};
    \node[anchor=north west, text width=\tw cm, align=left, font=\normalsize] at (0.3,3.45) {#3};
    \node[anchor=south east, font=\footnotesize, text=black!55] at (#1-0.2,0.15) {Draw a picture here.};
  \end{tikzpicture}}
% three cards across: \timelinecard[5.6]{...}{...}
\newcommand{\cuthere}{\noindent{\Large\ding{34}}\,{\small\itshape Cut out the cards. Add them to the class timeline.}\par\vspace{0.3em}}

% Build the key with:  pdflatex -jobname=<name>_KEY "\def\iskey{}\input{<name>}"
% Figures shared with the web page live in figures/Unit_X/<name>.tex and are
% pulled in with \fig{Unit_X/<name>}; the standalone package drops their own
% preamble, so their macros come from figures/preamble.tex, loaded here.
% Shared preamble for web figures.  Each figure is a standalone .tex that
% \input's this file; scripts/make-figures.sh turns it into an SVG.
% Handouts pull the same file with \fig{Unit_X/<name>}.
\usepackage[T1]{fontenc}
\usepackage{helvet}
\renewcommand{\familydefault}{\sfdefault}
\renewcommand{\ttdefault}{pcr}  % Courier: a real Type 1 typewriter font
\usepackage{amsmath}
\usepackage{tikz}
\usetikzlibrary{arrows.meta,positioning,calc,shapes.symbols,shapes.geometric}

% Who makes the rules is the point of most diagrams in this book, so the
% colors mean the same thing everywhere: blue is a person, orange is the
% computer, green is an answer, pink is the human brain.
\definecolor{person}{HTML}{2C6FAE}
\definecolor{machine}{HTML}{D35400}
\definecolor{result}{HTML}{1E8449}
\definecolor{brainpink}{HTML}{C2185B}
\definecolor{ink}{HTML}{2C3E50}

\tikzset{
  doodle/.style={line width=1.5pt, line cap=round, line join=round, draw=ink},
  box/.style={draw, rounded corners=5pt, line width=1.1pt, align=center,
              inner sep=7pt, minimum height=1.0cm, font=\small},
  personbox/.style={box, draw=person, fill=person!10},
  machinebox/.style={box, draw=machine, fill=machine!10},
  resultbox/.style={box, draw=result, fill=result!10},
  plainbox/.style={box, draw=ink!70, fill=ink!4},
  flow/.style={-{Stealth[length=8pt]}, line width=1.4pt, draw=ink},
}

% A person (blue) and a robot head (orange) as small icons.
% \personicon{x}{y}{scale}   \roboticon{x}{y}{scale}
\newcommand{\personicon}[3]{%
  \begin{scope}[shift={(#1,#2)}, scale=#3]
    \fill[person] (0,0.62) circle (0.28);
    \fill[person, rounded corners=6pt] (-0.42,-0.55) rectangle (0.42,0.28);
  \end{scope}}
\newcommand{\roboticon}[3]{%
  \begin{scope}[shift={(#1,#2)}, scale=#3]
    \draw[machine, line width=1.2pt] (0,0.5) -- (0,0.72);
    \fill[machine] (0,0.78) circle (0.08);
    \draw[machine, line width=1.4pt, fill=machine!12, rounded corners=4pt] (-0.5,-0.45) rectangle (0.5,0.5);
    \fill[ink] (-0.2,0.12) circle (0.08);
    \fill[ink] (0.2,0.12) circle (0.08);
    \draw[ink, line width=1pt] (-0.22,-0.22) -- (0.22,-0.22);
  \end{scope}}

\usepackage[margin=0.75in]{geometry}
\usepackage{array}
\usepackage{enumitem}
\usepackage{amsmath}
\usepackage{cancel}
\usepackage[helvet]{sfmath}  % math in the same sans font as the text
\usepackage{standalone}
\renewcommand{\ttdefault}{pcr}  % Courier: a real Type 1 typewriter font
\pagestyle{empty}
\setlength{\parindent}{0pt}
\setlist[enumerate,1]{leftmargin=1.5em,itemsep=1em}
\setlist[enumerate,2]{label=(\alph*),leftmargin=2em,itemsep=2.2em}

\newif\ifkey
\ifdefined\iskey \keytrue \fi
\colorlet{keyred}{red!75!black}
\newcommand{\ans}[1]{\ifkey\leavevmode{\color{keyred}\let\frac\dfrac #1}\fi}  % full-size fractions in answers
% vertical work space on the student copy, less on the key
\newcommand{\work}[1]{\ifkey\vspace{0.6em}\else\vspace{#1}\fi}
\newcommand{\fig}[1]{\input{../../figures/#1}}
% Title line: \lessontitle{Quartiles, IQR, and Box Plots}
% A long title gets the Name line on its own line above it.
\newlength{\titlewidth}
\newcommand{\lessontitle}[1]{%
  \settowidth{\titlewidth}{\Large\bfseries #1\ifkey\ KEY\fi}%
  \ifdim\titlewidth>4.3in
    \hfill Name \underline{\hspace{2in}}\par\vspace{0.5em}
    {\Large\bfseries #1\ifkey\ \textcolor{keyred}{KEY}\fi}\par\vspace{0.6em}
  \else
    {\Large\bfseries #1\ifkey\ \textcolor{keyred}{KEY}\fi}\hfill Name \underline{\hspace{2in}}\par\vspace{0.6em}
  \fi}
\newcommand{\partner}{\par\vspace{0.4em}Work with a partner. Compare answers when you finish each problem.\par}

% ---------- shared pieces added for AI and You ----------
% no hyphenated words anywhere: students read in a second language
\hyphenpenalty=10000 \exhyphenpenalty=10000
\usepackage{pifont}   % \ding{34} scissors
\newcommand{\blank}[1]{\underline{\hspace{#1}}}
% ragged table columns: no hyphenated words for students reading in a second language
\newcolumntype{R}[1]{>{\raggedright\arraybackslash}p{#1}}
% An answer written on a blank line: \ansline{width}{answer}
\newcommand{\ansline}[2]{\ifkey\underline{\makebox[#1][l]{\,\textcolor{keyred}{#2}}}\else\blank{#1}\fi}
% Unit 1 class timeline: a card to cut out and post on the wall.
%   \timelinecard{1950}{Alan Turing asks ...}
\newcommand{\timelinecard}[3][8.3]{%
  \begin{tikzpicture}
    \pgfmathsetmacro{\tw}{#1-0.8}
    \draw[dashed, line width=0.8pt, rounded corners=3pt] (0,0) rectangle (#1,4.7);
    \node[anchor=north west, font=\huge\bfseries] at (0.3,4.45) {#2};
    \node[anchor=north west, text width=\tw cm, align=left, font=\normalsize] at (0.3,3.45) {#3};
    \node[anchor=south east, font=\footnotesize, text=black!55] at (#1-0.2,0.15) {Draw a picture here.};
  \end{tikzpicture}}
% three cards across: \timelinecard[5.6]{...}{...}
\newcommand{\cuthere}{\noindent{\Large\ding{34}}\,{\small\itshape Cut out the cards. Add them to the class timeline.}\par\vspace{0.3em}}

% Build the key with:  pdflatex -jobname=<name>_KEY "\def\iskey{}\input{<name>}"
% Figures shared with the web page live in figures/Unit_X/<name>.tex and are
% pulled in with \fig{Unit_X/<name>}; the standalone package drops their own
% preamble, so their macros come from figures/preamble.tex, loaded here.
% Shared preamble for web figures.  Each figure is a standalone .tex that
% \input's this file; scripts/make-figures.sh turns it into an SVG.
% Handouts pull the same file with \fig{Unit_X/<name>}.
\usepackage[T1]{fontenc}
\usepackage{helvet}
\renewcommand{\familydefault}{\sfdefault}
\renewcommand{\ttdefault}{pcr}  % Courier: a real Type 1 typewriter font
\usepackage{amsmath}
\usepackage{tikz}
\usetikzlibrary{arrows.meta,positioning,calc,shapes.symbols,shapes.geometric}

% Who makes the rules is the point of most diagrams in this book, so the
% colors mean the same thing everywhere: blue is a person, orange is the
% computer, green is an answer, pink is the human brain.
\definecolor{person}{HTML}{2C6FAE}
\definecolor{machine}{HTML}{D35400}
\definecolor{result}{HTML}{1E8449}
\definecolor{brainpink}{HTML}{C2185B}
\definecolor{ink}{HTML}{2C3E50}

\tikzset{
  doodle/.style={line width=1.5pt, line cap=round, line join=round, draw=ink},
  box/.style={draw, rounded corners=5pt, line width=1.1pt, align=center,
              inner sep=7pt, minimum height=1.0cm, font=\small},
  personbox/.style={box, draw=person, fill=person!10},
  machinebox/.style={box, draw=machine, fill=machine!10},
  resultbox/.style={box, draw=result, fill=result!10},
  plainbox/.style={box, draw=ink!70, fill=ink!4},
  flow/.style={-{Stealth[length=8pt]}, line width=1.4pt, draw=ink},
}

% A person (blue) and a robot head (orange) as small icons.
% \personicon{x}{y}{scale}   \roboticon{x}{y}{scale}
\newcommand{\personicon}[3]{%
  \begin{scope}[shift={(#1,#2)}, scale=#3]
    \fill[person] (0,0.62) circle (0.28);
    \fill[person, rounded corners=6pt] (-0.42,-0.55) rectangle (0.42,0.28);
  \end{scope}}
\newcommand{\roboticon}[3]{%
  \begin{scope}[shift={(#1,#2)}, scale=#3]
    \draw[machine, line width=1.2pt] (0,0.5) -- (0,0.72);
    \fill[machine] (0,0.78) circle (0.08);
    \draw[machine, line width=1.4pt, fill=machine!12, rounded corners=4pt] (-0.5,-0.45) rectangle (0.5,0.5);
    \fill[ink] (-0.2,0.12) circle (0.08);
    \fill[ink] (0.2,0.12) circle (0.08);
    \draw[ink, line width=1pt] (-0.22,-0.22) -- (0.22,-0.22);
  \end{scope}}

\usepackage[margin=0.75in]{geometry}
\usepackage{array}
\usepackage{enumitem}
\usepackage{amsmath}
\usepackage{cancel}
\usepackage[helvet]{sfmath}  % math in the same sans font as the text
\usepackage{standalone}
\renewcommand{\ttdefault}{pcr}  % Courier: a real Type 1 typewriter font
\pagestyle{empty}
\setlength{\parindent}{0pt}
\setlist[enumerate,1]{leftmargin=1.5em,itemsep=1em}
\setlist[enumerate,2]{label=(\alph*),leftmargin=2em,itemsep=2.2em}

\newif\ifkey
\ifdefined\iskey \keytrue \fi
\colorlet{keyred}{red!75!black}
\newcommand{\ans}[1]{\ifkey\leavevmode{\color{keyred}\let\frac\dfrac #1}\fi}  % full-size fractions in answers
% vertical work space on the student copy, less on the key
\newcommand{\work}[1]{\ifkey\vspace{0.6em}\else\vspace{#1}\fi}
\newcommand{\fig}[1]{\input{../../figures/#1}}
% Title line: \lessontitle{Quartiles, IQR, and Box Plots}
% A long title gets the Name line on its own line above it.
\newlength{\titlewidth}
\newcommand{\lessontitle}[1]{%
  \settowidth{\titlewidth}{\Large\bfseries #1\ifkey\ KEY\fi}%
  \ifdim\titlewidth>4.3in
    \hfill Name \underline{\hspace{2in}}\par\vspace{0.5em}
    {\Large\bfseries #1\ifkey\ \textcolor{keyred}{KEY}\fi}\par\vspace{0.6em}
  \else
    {\Large\bfseries #1\ifkey\ \textcolor{keyred}{KEY}\fi}\hfill Name \underline{\hspace{2in}}\par\vspace{0.6em}
  \fi}
\newcommand{\partner}{\par\vspace{0.4em}Work with a partner. Compare answers when you finish each problem.\par}

% ---------- shared pieces added for AI and You ----------
% no hyphenated words anywhere: students read in a second language
\hyphenpenalty=10000 \exhyphenpenalty=10000
\usepackage{pifont}   % \ding{34} scissors
\newcommand{\blank}[1]{\underline{\hspace{#1}}}
% ragged table columns: no hyphenated words for students reading in a second language
\newcolumntype{R}[1]{>{\raggedright\arraybackslash}p{#1}}
% An answer written on a blank line: \ansline{width}{answer}
\newcommand{\ansline}[2]{\ifkey\underline{\makebox[#1][l]{\,\textcolor{keyred}{#2}}}\else\blank{#1}\fi}
% Unit 1 class timeline: a card to cut out and post on the wall.
%   \timelinecard{1950}{Alan Turing asks ...}
\newcommand{\timelinecard}[3][8.3]{%
  \begin{tikzpicture}
    \pgfmathsetmacro{\tw}{#1-0.8}
    \draw[dashed, line width=0.8pt, rounded corners=3pt] (0,0) rectangle (#1,4.7);
    \node[anchor=north west, font=\huge\bfseries] at (0.3,4.45) {#2};
    \node[anchor=north west, text width=\tw cm, align=left, font=\normalsize] at (0.3,3.45) {#3};
    \node[anchor=south east, font=\footnotesize, text=black!55] at (#1-0.2,0.15) {Draw a picture here.};
  \end{tikzpicture}}
% three cards across: \timelinecard[5.6]{...}{...}
\newcommand{\cuthere}{\noindent{\Large\ding{34}}\,{\small\itshape Cut out the cards. Add them to the class timeline.}\par\vspace{0.3em}}

\renewcommand{\arraystretch}{1.5}

\begin{document}

\lessontitle{Simple Rules, Complex Behavior}
\partner

\wordbank[You can use a word more than once.]{\wb{2}\wb{30}\wb{a pattern that repeats}\wb{chaos}\wb{emergence}\wb{the wave}\wb{triangles}}

{\large\textbf{1. The wave.}} Was anyone in charge of the wave? \ansline{0.8in}{no}

\vspace{0.5em}
Did the person you watched move down the row? \ansline{0.8in}{no} \quad So what moved? \ansline{1.6in}{the wave}

\vspace{0.6em}
When many simple parts make something bigger that nobody planned, it's called \ansline{1.4in}{emergence}.

\vspace{0.8em}
{\large\textbf{2. Make a few rows.}} Look at a square and its two neighbors in the row above. Find the same 3 colors on the card. Color the square below. (Outside the grid, everything is white.)

\begin{center}\resizebox{0.95\linewidth}{!}{\fig{Unit_2/2.1_rule90_card}}\end{center}

\vspace{0.1em}
\begin{center}% generated by scripts/make-rule30.py: Rule 90, 5 rows, 13 cells
\begin{tikzpicture}
\fill[black] (4.680,3.120) rectangle ++(0.78,0.78);
\ifkey
\fill[keyred!70] (3.900,2.340) rectangle ++(0.78,0.78);
\fill[keyred!70] (5.460,2.340) rectangle ++(0.78,0.78);
\fill[keyred!70] (3.120,1.560) rectangle ++(0.78,0.78);
\fill[keyred!70] (6.240,1.560) rectangle ++(0.78,0.78);
\fill[keyred!70] (2.340,0.780) rectangle ++(0.78,0.78);
\fill[keyred!70] (3.900,0.780) rectangle ++(0.78,0.78);
\fill[keyred!70] (5.460,0.780) rectangle ++(0.78,0.78);
\fill[keyred!70] (7.020,0.780) rectangle ++(0.78,0.78);
\fill[keyred!70] (1.560,0.000) rectangle ++(0.78,0.78);
\fill[keyred!70] (7.800,0.000) rectangle ++(0.78,0.78);
\fi
\draw[black!55, line width=0.4pt] (0,0) grid[step=0.78] (10.140,3.900);
\node[anchor=east, font=\small] at (-0.1,3.510) {1};
\node[anchor=east, font=\small] at (-0.1,2.730) {2};
\node[anchor=east, font=\small] at (-0.1,1.950) {3};
\node[anchor=east, font=\small] at (-0.1,1.170) {4};
\node[anchor=east, font=\small] at (-0.1,0.390) {5};
\end{tikzpicture}
\end{center}

\vspace{0.2em}
What shape is starting to appear? \ansline{1.4in}{triangles}

\vspace{0.8em}
{\large\textbf{3. One small change.}} How many answers are different on the Rule 30 card? \ansline{0.6in}{2}

\vspace{0.5em}
Rule 90 makes \ansline{2.0in}{a pattern that repeats}. \quad Rule 30 makes \ansline{1.0in}{chaos}.

\vspace{0.5em}
The cone snail's shell looks like Rule \ansline{0.6in}{30}.

\newpage

\wordbank{\wb{2 or 3}\wb{3}\wb{backward}\wb{the pattern}}

{\large\textbf{4. Explore the rules.}} Draw each rule. Circle what it looks like.

\vspace{0.3em}
\begin{tabular}{|c|R{5.6in}|}
\hline
\textbf{Rule} & \textbf{What does it look like?} \\ \hline
110 & never changes \qquad a pattern that repeats \qquad chaos \qquad complex \ \ans{(complex)} \\ \hline
4 & never changes \qquad a pattern that repeats \qquad chaos \qquad complex \ \ans{(never changes)} \\ \hline
150 & never changes \qquad a pattern that repeats \qquad chaos \qquad complex \ \ans{(repeats)} \\ \hline
30 & never changes \qquad a pattern that repeats \qquad chaos \qquad complex \ \ans{(chaos)} \\ \hline
\end{tabular}

\vspace{0.5em}
\textbf{My own rule.} My rule number is \ansline{1.1in}{Answers vary.} It looks like \ansline{2.3in}{}

\vspace{0.6em}
\textbf{Traffic (Rule 184).} Every car moves forward. Which way do the jams move? \ansline{1.6in}{backward}

\vspace{0.8em}
{\large\textbf{5. Simple rules everywhere.}}

\vspace{0.4em}
Without ``fly the same way,'' the birds \ansline{4.0in}{stay in a crowd, but each points its own way.}

\vspace{0.5em}
In the Game of Life, did any single cell move? \ansline{0.6in}{no} \quad So what moved? \ansline{1.8in}{the pattern}

\vspace{0.8em}
{\large\textbf{6. Talk about it.}} Can simple parts make something smart?

\vspace{0.4em}
I think simple rules \textbf{can} / \textbf{cannot} make something smart, because \ansline{2.5in}{Answers vary.}

\vspace{0.8em}
\hrule
\vspace{0.6em}
{\large\textbf{Go Further (Day 2).}} Be Rule 30 as a class. Color each new row the class shows.

\vspace{0.2em}
\begin{center}% generated by scripts/make-rule30.py: Rule 30, 10 rows, 21 cells
\begin{tikzpicture}
\fill[black] (7.800,7.020) rectangle ++(0.78,0.78);
\ifkey
\fill[keyred!70] (7.020,6.240) rectangle ++(0.78,0.78);
\fill[keyred!70] (7.800,6.240) rectangle ++(0.78,0.78);
\fill[keyred!70] (8.580,6.240) rectangle ++(0.78,0.78);
\fill[keyred!70] (6.240,5.460) rectangle ++(0.78,0.78);
\fill[keyred!70] (7.020,5.460) rectangle ++(0.78,0.78);
\fill[keyred!70] (9.360,5.460) rectangle ++(0.78,0.78);
\fill[keyred!70] (5.460,4.680) rectangle ++(0.78,0.78);
\fill[keyred!70] (6.240,4.680) rectangle ++(0.78,0.78);
\fill[keyred!70] (7.800,4.680) rectangle ++(0.78,0.78);
\fill[keyred!70] (8.580,4.680) rectangle ++(0.78,0.78);
\fill[keyred!70] (9.360,4.680) rectangle ++(0.78,0.78);
\fill[keyred!70] (10.140,4.680) rectangle ++(0.78,0.78);
\fill[keyred!70] (4.680,3.900) rectangle ++(0.78,0.78);
\fill[keyred!70] (5.460,3.900) rectangle ++(0.78,0.78);
\fill[keyred!70] (7.800,3.900) rectangle ++(0.78,0.78);
\fill[keyred!70] (10.920,3.900) rectangle ++(0.78,0.78);
\fill[keyred!70] (3.900,3.120) rectangle ++(0.78,0.78);
\fill[keyred!70] (4.680,3.120) rectangle ++(0.78,0.78);
\fill[keyred!70] (6.240,3.120) rectangle ++(0.78,0.78);
\fill[keyred!70] (7.020,3.120) rectangle ++(0.78,0.78);
\fill[keyred!70] (7.800,3.120) rectangle ++(0.78,0.78);
\fill[keyred!70] (8.580,3.120) rectangle ++(0.78,0.78);
\fill[keyred!70] (10.140,3.120) rectangle ++(0.78,0.78);
\fill[keyred!70] (10.920,3.120) rectangle ++(0.78,0.78);
\fill[keyred!70] (11.700,3.120) rectangle ++(0.78,0.78);
\fill[keyred!70] (3.120,2.340) rectangle ++(0.78,0.78);
\fill[keyred!70] (3.900,2.340) rectangle ++(0.78,0.78);
\fill[keyred!70] (6.240,2.340) rectangle ++(0.78,0.78);
\fill[keyred!70] (10.140,2.340) rectangle ++(0.78,0.78);
\fill[keyred!70] (12.480,2.340) rectangle ++(0.78,0.78);
\fill[keyred!70] (2.340,1.560) rectangle ++(0.78,0.78);
\fill[keyred!70] (3.120,1.560) rectangle ++(0.78,0.78);
\fill[keyred!70] (4.680,1.560) rectangle ++(0.78,0.78);
\fill[keyred!70] (5.460,1.560) rectangle ++(0.78,0.78);
\fill[keyred!70] (6.240,1.560) rectangle ++(0.78,0.78);
\fill[keyred!70] (7.020,1.560) rectangle ++(0.78,0.78);
\fill[keyred!70] (9.360,1.560) rectangle ++(0.78,0.78);
\fill[keyred!70] (10.140,1.560) rectangle ++(0.78,0.78);
\fill[keyred!70] (10.920,1.560) rectangle ++(0.78,0.78);
\fill[keyred!70] (11.700,1.560) rectangle ++(0.78,0.78);
\fill[keyred!70] (12.480,1.560) rectangle ++(0.78,0.78);
\fill[keyred!70] (13.260,1.560) rectangle ++(0.78,0.78);
\fill[keyred!70] (1.560,0.780) rectangle ++(0.78,0.78);
\fill[keyred!70] (2.340,0.780) rectangle ++(0.78,0.78);
\fill[keyred!70] (4.680,0.780) rectangle ++(0.78,0.78);
\fill[keyred!70] (7.800,0.780) rectangle ++(0.78,0.78);
\fill[keyred!70] (8.580,0.780) rectangle ++(0.78,0.78);
\fill[keyred!70] (9.360,0.780) rectangle ++(0.78,0.78);
\fill[keyred!70] (14.040,0.780) rectangle ++(0.78,0.78);
\fill[keyred!70] (0.780,0.000) rectangle ++(0.78,0.78);
\fill[keyred!70] (1.560,0.000) rectangle ++(0.78,0.78);
\fill[keyred!70] (3.120,0.000) rectangle ++(0.78,0.78);
\fill[keyred!70] (3.900,0.000) rectangle ++(0.78,0.78);
\fill[keyred!70] (4.680,0.000) rectangle ++(0.78,0.78);
\fill[keyred!70] (5.460,0.000) rectangle ++(0.78,0.78);
\fill[keyred!70] (7.020,0.000) rectangle ++(0.78,0.78);
\fill[keyred!70] (7.800,0.000) rectangle ++(0.78,0.78);
\fill[keyred!70] (10.140,0.000) rectangle ++(0.78,0.78);
\fill[keyred!70] (13.260,0.000) rectangle ++(0.78,0.78);
\fill[keyred!70] (14.040,0.000) rectangle ++(0.78,0.78);
\fill[keyred!70] (14.820,0.000) rectangle ++(0.78,0.78);
\fi
\draw[black!55, line width=0.4pt] (0,0) grid[step=0.78] (16.380,7.800);
\node[anchor=east, font=\small] at (-0.1,7.410) {1};
\node[anchor=east, font=\small] at (-0.1,6.630) {2};
\node[anchor=east, font=\small] at (-0.1,5.850) {3};
\node[anchor=east, font=\small] at (-0.1,5.070) {4};
\node[anchor=east, font=\small] at (-0.1,4.290) {5};
\node[anchor=east, font=\small] at (-0.1,3.510) {6};
\node[anchor=east, font=\small] at (-0.1,2.730) {7};
\node[anchor=east, font=\small] at (-0.1,1.950) {8};
\node[anchor=east, font=\small] at (-0.1,1.170) {9};
\node[anchor=east, font=\small] at (-0.1,0.390) {10};
\end{tikzpicture}
\end{center}

\vspace{0.4em}
\textbf{Game of Life.} A black cell with \ansline{0.9in}{2 or 3} black neighbors stays black. A white cell with exactly \ansline{0.6in}{3} black neighbors turns black.

\end{document}
